\section{Qualité et préparation des données avec Spark}
\begin{itemize}
  \item Présenter le chargement du DataFrame.
  \item Afficher et commenter \texttt{printSchema()}.
  \item Contrôler les NULL, doublons et types.
  \item Présenter les règles de nettoyage réellement appliquées et leur justification.
\end{itemize}
\textbf{[TODO — team:data-eng]}

\begin{table}[H]
\centering
\caption{Contrôles qualité appliqués}
\begin{tabular}{llll}
\toprule
Contrôle & Résultat & Action réalisée & Justification \\
\midrule
Valeurs NULL & & & \\
Doublons & & & \\
Types & & & \\
Valeurs incohérentes & & & \\
\bottomrule
\end{tabular}
\end{table}
